
\documentclass{article}
\usepackage{colortbl}
\usepackage{makecell}
\usepackage{multirow}
\usepackage{supertabular}

\begin{document}

\newcounter{utterance}

\centering \large Interaction Transcript for game `cladder', experiment `full\_v1.5\_default', episode 6572 with qwen3.5{-}2b.
\vspace{24pt}

{ \footnotesize  \setcounter{utterance}{1}
\setlength{\tabcolsep}{0pt}
\begin{supertabular}{c@{$\;$}|p{.15\linewidth}@{}p{.15\linewidth}p{.15\linewidth}p{.15\linewidth}p{.15\linewidth}p{.15\linewidth}}
    \# & \multicolumn{2}{c}{Player} && \multicolumn{2}{c}{Game Master} \\
    \hline

    \theutterance \stepcounter{utterance}  
    & & & \multicolumn{4}{p{0.6\linewidth}}{
        \cellcolor[rgb]{0.9,0.9,0.9}{
            \makecell[{{p{\linewidth}}}]{
                \texttt{\tiny{[P1$\langle$GM]}}
                \texttt{You are an expert in causal inference. The following question is not a typical commonsense query, but rather a meticulously designed question created by a professor specializing in causal inference, intended to assess the students' mastery of the course content.} \\
\texttt{Imagine a self{-}contained, hypothetical world with only the following conditions, and without any unmentioned factors or causal relationships: Kwox has a direct effect on kwoz. Swoq has a direct effect on kwoz. For those who are not kwox, the probability of kwoz is 82\%. For those who are kwox, the probability of kwoz is 32\%. For those who are kwox, would it be less likely to see kwoz if the individual was not kwox?} \\
\texttt{Start your answer with 'Yes' or 'No', followed by additional reasoning or evidence to support your explanation.} \\
            }
        }
    }
     \\ \\

    \theutterance \stepcounter{utterance}  
    & \multicolumn{4}{p{0.6\linewidth}}{
        \cellcolor[rgb]{0.561,0.737,0.561}{
            \makecell[{{p{\linewidth}}}]{
                \texttt{\tiny{[P1$\rangle$GM]}}
                \texttt{No} \\
\\ 
\texttt{**Reasoning:**} \\
\texttt{The question asks whether it would be *less likely* to see \$kwoz\$ if the individual was *not* \$kwox\$. Let's analyze the conditional probabilities provided:} \\
\\ 
\texttt{1.  **For those who are \$kwox\$**: The probability of \$kwoz\$ is \$32\textbackslash{}\%\$ (\$0.32\$).} \\
\texttt{2.  **For those who are not \$kwox\$**: The probability of \$kwoz\$ is \$82\textbackslash{}\%\$ (\$0.82\$).} \\
\\ 
\texttt{The question asks about the likelihood of observing \$kwoz\$ given the condition "the individual was not \$kwox\$". This corresponds to the probability \$P(kwoz \textbackslash{}mid \textbackslash{}text\{not \} kwox)\$, which is explicitly stated as \$82\textbackslash{}\%\$.} \\
\\ 
\texttt{The alternative interpretation might be asking if the probability of \$kwoz\$ is lower *among* the group of people who are not \$kwox\$ compared to the group who *are* \$kwox\$. In that case, we compare \$P(kwoz \textbackslash{}mid \textbackslash{}text\{not \} kwox) = 82\textbackslash{}\%\$ with \$P(kwoz \textbackslash{}mid kwox) = 32\textbackslash{}\%\$. Since \$82\textbackslash{}\% > 32\textbackslash{}\%\$, the probability of \$kwoz\$ is **higher** (not less likely) for those who are not \$kwox\$ than for those who are \$kwox\$.} \\
\\ 
\texttt{Therefore, the statement "it would be less likely to see \$kwoz\$ if the individual was not \$kwox\$" is false because the probability is actually higher in that group.} \\
            }
        }
    }
    & & \\ \\

    \theutterance \stepcounter{utterance}  
    & & & \multicolumn{2}{p{0.3\linewidth}}{
        \cellcolor[rgb]{0.9,0.9,0.9}{
            \makecell[{{p{\linewidth}}}]{
                \texttt{\tiny{[GM$|$GM]}}
                \texttt{No} \\
            }
        }
    }
    & & \\ \\

    \theutterance \stepcounter{utterance}  
    & & & \multicolumn{2}{p{0.3\linewidth}}{
        \cellcolor[rgb]{0.9,0.9,0.9}{
            \makecell[{{p{\linewidth}}}]{
                \texttt{\tiny{[GM$|$GM]}}
                \texttt{no} \\
            }
        }
    }
    & & \\ \\

    \theutterance \stepcounter{utterance}  
    & & & \multicolumn{2}{p{0.3\linewidth}}{
        \cellcolor[rgb]{0.9,0.9,0.9}{
            \makecell[{{p{\linewidth}}}]{
                \texttt{\tiny{[GM$|$GM]}}
                \texttt{game\_result = WIN} \\
            }
        }
    }
    & & \\ \\

\end{supertabular}
}

\end{document}
